%% data-access.tex
%% Shared database access and iteration helpers for ttrpg-core.

%% \get{key} — expands to the current entry's value for that key.
%% Nested objects: \get{stats@hp}   Arrays: \get{tags} returns comma-separated values.
%% For specific array indices: \get{tags@1}, \get{tags@count} still work.
\newcount\ttrpg@arrayindex
\newcommand{\getArrayContents}[2][,]{%
    \ttrpg@arrayindex=1\relax%
    \loop%
    \ifnum\ttrpg@arrayindex<\csname db@#2@count\endcsname\relax%
    % Use \unexpanded to protect the delimiter from premature expansion
    \csname db@#2@\the\ttrpg@arrayindex\endcsname\unexpanded{#1}%
    \advance\ttrpg@arrayindex by 1\relax%
    \repeat%
    \ifnum\ttrpg@arrayindex=\csname db@#2@count\endcsname\relax%
        \csname db@#2@\the\ttrpg@arrayindex\endcsname%
    \fi%
}
\newcommand{\get}[1]{%
    \ifcsname db@#1@count\endcsname
        \getArrayContents[,]{#1}%
    \else
        \csname db@#1\endcsname
    \fi
}

%% \ifhas{key}{code} — runs code only if key was present in the JSON.
%% Existence is scoped to the current entry (\currentfile) so optional
%% fields from prior entries do not leak forward.
\newcommand{\ifhas}[2]{%
    \ifcsname db@#1@exists\endcsname
        \edef\ttrpg@existsmarker{\csname db@#1@exists\endcsname}%
        \edef\ttrpg@currententry{\currentfile}%
        \ifdefstrequal{\ttrpg@existsmarker}{\ttrpg@currententry}{#2}{}%
    \fi%
}

%% \ifhaselse{key}{then}{else} — runs then/else based on key presence in current entry.
\newcommand{\ifhaselse}[3]{%
    \ifcsname db@#1@exists\endcsname
        \edef\ttrpg@existsmarker{\csname db@#1@exists\endcsname}%
        \edef\ttrpg@currententry{\currentfile}%
        \ifdefstrequal{\ttrpg@existsmarker}{\ttrpg@currententry}{#2}{#3}%
    \else
        #3%
    \fi%
}

%% \getOrDefault{key}{fallback} — expands to key value if present, else fallback.
\newcommand{\getOrDefault}[2]{%
    \ifcsname db@#1@exists\endcsname
        \csname db@#1\endcsname
    \else
        #2%
    \fi
}

%% Lowercase alias for template convenience.
\let\getordefault\getOrDefault

%% Key-iteration helpers for object fields (e.g. class_lists).

%% Lookup helpers used by \forEachKeyValue.
\newcommand{\dbkeyname}[2]{%
    \csname db@#1@keys@#2\endcsname%
}

\newcommand{\dbkeysafe}[2]{%
    \csname db@#1@keys@#2@safe\endcsname%
}

\newcommand{\dbkeycount}[1]{%
    \ifcsname db@#1@keys@count\endcsname
        \csname db@#1@keys@count\endcsname
    \else
        0%
    \fi
}

\newcommand{\ifKeyValueType}[3]{\ifdefstring{\KeyValueType}{#1}{#2}{#3}}

\newcount\ttrpg@keyindex
\newcount\ttrpg@keycount

\newcommand{\forEachKeyValue}[2]{%
    \ttrpg@keyindex=1\relax
    \ttrpg@keycount=\numexpr\dbkeycount{#1}\relax
    \edef\KeyCount{\the\ttrpg@keycount}%
    \advance\ttrpg@keycount by 1\relax
    \loop
    \ifnum\ttrpg@keyindex<\ttrpg@keycount
    \edef\KeyIndex{\the\ttrpg@keyindex}%
    \edef\KeyName{\dbkeyname{#1}{\KeyIndex}}%
    \edef\KeySafe{\dbkeysafe{#1}{\KeyIndex}}%
    \def\KeyValue{\csname db@#1@\KeySafe\endcsname}%
    \ifcsname db@#1@\KeySafe @type\endcsname%
        \edef\KeyValueType{\csname db@#1@\KeySafe @type\endcsname}%
    \else%
        \edef\KeyValueType{unknown}%
    \fi%
    #2%
    \advance\ttrpg@keyindex by 1\relax
    \repeat
}

%% \CRtoXP[system]{cr} — returns XP value for a CR.
%% Default system is PF1E. Unknown CR/system returns ?.
\newcommand{\CRtoXP}[2][pf1e]{%
    \directlua{tex.sprint(ttrpg_xp.cr_to_xp("\luaescapestring{#1}", "\luaescapestring{#2}"))}%
}
